\documentclass[11pt,a4paper]{article}
\usepackage[utf8]{inputenc}
\usepackage{amsmath}
\usepackage{amssymb}
\usepackage{booktabs}
\usepackage{geometry}
\usepackage{microtype}
\usepackage{hyperref}
\geometry{margin=1in}

\title{\textbf{Technical Supplement \& Parameter Index:\\
		Multi-Scale Chameleon Scalar Fields and Anisotropic Speed of Light Profiles}}
\author{\textbf{Phillip A. Longo} \\ ORCID: 0009-0004-3537-6478}
\date{\today}

\begin{document}
	
	\maketitle
	
	\begin{abstract}
		This technical supplement details the physical parameters, field equations, and boundary conditions evaluated across the 10-module integration pipeline and its associated numerical simulations. This index documents the parameter configurations utilized in the accompanying preprint repository, which evaluates a six-dimensional semiclassical action functional over an extended configuration manifold to address the local cosmological expansion discrepancy via a transient early dark energy contribution while satisfying local terrestrial screening constraints.
	\end{abstract}
	\section{Pipeline Framework Structure}
	The technical supplement and its corresponding repository code operate via a 10-module synchronized calculation pipeline, structured as follows:
	\begin{enumerate}
		\item \textbf{Vacuum Baseline Array:} Initializes standard cosmological parameters and energy benchmarks.
		\item \textbf{Sound Horizon Root-Solver:} Automatically derives the geometric Early Dark Energy tracking scale factor.
		\item \textbf{Torsion Scale Converter:} Maps high-energy cutoff thresholds down to absolute length dimensions.
		\item \textbf{Velocity Dispersion Filter:} Linearizes metric perturbations to evaluate deep-space wave propagation.
		\item \textbf{Chameleon Screening Module:} Pins local scalar gradients within massive material boundaries.
		\item \textbf{Lapse Decoherence Tracking:} Monitors temporal tracking phase variations under strong-field boundaries.
		\item \textbf{Quantum Damping Profile:} Regulates non-unitary state suppression as the gravitational lapse vanishes.
		\item \textbf{Cosmological Delay Integral:} Evaluates high-redshift multi-messenger lookback travel times.
		\item \textbf{Substrate Shielding Engine:} Simulates phonon collision minimization in condensed matter.
		\item \textbf{SQUID Transport Mapping:} Transforms phase-slip flux directly into readable voltage modulation curves.
	\end{enumerate}
	
	\section{Simulation Constants and Boundary Baselines}
	The numerical simulation environment fixes its baseline variables using the cosmological and environmental constraints detailed in Table 1:
	\begin{table}[h!]
		\centering
		\small
		\begin{tabular}{llp{1.4in}p{2.3in}}
			\toprule
			\textbf{Parameter Anchor} & \textbf{Sym.} & \textbf{Calibrated Value} & \textbf{Physical Definition / Boundary} \\
			\midrule
			Absolute vacuum speed of light & $c$ & $299792458.0 \text{ m/s}$ & Fixed laboratory baseline \\
			Early Dark Energy Density Peak & $f_{\text{EDE}}$ & \textbf{0.0589} & Peak fraction at critical redshift $z_c \approx 3500$ \\
			Sound Horizon Compression Factor & $\Delta r_s^*$ & \textbf{9.9497\%} & Comoving sound horizon scale modification \\
			Early Reconciled Expansion Anchor & $H_{0,\text{rec}}$ & \textbf{67.95 km/s/Mpc} & Reconciled CMB acoustic horizon baseline \\
			Maximum local velocity rift & $\Delta c_{\text{max}}$ & $1.7097 \text{ m/s}$ & Unscreened deep-space anisotropy peak \\
			Minimum anisotropic velocity bound & $c_{\text{min}}$ & $299792456.2903 \text{ m/s}$ & Axis parameter under maximal deceleration \\
			Effective field theory cutoff scale & $\Lambda_{\text{EFT}}$ & $1.0000 \times 10^{-11} \text{ J}$ & Short-range ultraviolet energy regulator \\
			Reconciled Local Hubble constant & $H_0$ & \textbf{74.39 km/s/Mpc} & Fully resolved local expansion rate \\
			\bottomrule
		\end{tabular}
		\caption{Reconciled simulation constraints matching repository software diagnostic outputs.}
	\end{table}
	
	\section{Field Equations and Local Screening Constraints}
	The directional velocity modification profile is mapped relative to an external scalar gradient via the spatial metric perturbation tensor projection:
	\begin{equation}
		c_{\text{eff}}(\hat{\mathbf{k}}) = c \left[ 1 - \frac{\alpha_0^2 \lambda_{\text{torsion}}^2 (1 - \Phi^2)^2}{2} |\nabla\Phi \times \hat{\mathbf{k}}|^2 \right]
	\end{equation}
	
	In high-density material environments ($\rho_b \gg 1.0$), the effective mass of the scalar field increases dynamically, pinning the tracking field variable to its boundary limits ($\Phi \rightarrow \pm 1$). This suppresses spatial gradients ($\nabla\Phi \rightarrow 0$), satisfying local experimental constraints down to the $10^{-17}$ level within cryogenic cavities to mask macroscopic vacuum anisotropies:
	\begin{equation}
		\mathcal{F}_{\text{screening}}(\rho_b) \propto e^{-\alpha_0^2 \cdot 550} \Longrightarrow c_{\text{eff}}(\theta) \longrightarrow c
	\end{equation}
	
	Under high space-time curvature boundaries near the gravitational lapse limit ($N \rightarrow 0$), the temporal component of the field's kinetic tracking term undergoes a rapid phase divergence to preserve high-energy unitarity, causing instantaneous state decoherence:
	\begin{equation}
		\Gamma_{nn'}(N) = \frac{\alpha_0^2 \omega^2 \tau_J^2}{3 N^2} \left( \frac{\hbar \omega}{\Lambda_{\text{EFT}}} \right)^3 \left( \frac{\partial \Phi}{\partial t} \right)^2 \Delta n^2 \longrightarrow +\infty
	\end{equation}
	where $\tau_J \equiv \hbar / (2e I_0 R)$ is the characteristic Josephson relaxation time, $\omega$ is the operational transition frequency, and $\Delta n$ represents the quantum state difference. This positive divergence enforces stable non-unitary damping as the localized metric configurations near the horizon limit.
	
	\section{Cosmological Timing Delays and Coherence Shielding}
	Integrating photon lookback travel times across high-redshift boundaries ($z_{\text{source}} = 2.0$) relative to a gravitational wave counterpart yields an explicit arrival time separation at Earth due to the underlying field coupling variations. Direct analytical separation isolation yields:
	\begin{equation}
		\Delta t_{\text{arrival}} = \int_{0}^{z_{\text{source}}} \frac{dz}{(1+z)H(z)} \left[ \frac{1}{v_{\text{photon}}(z)} - \frac{1}{v_{\text{gw}}(z)} \right] = \mathbf{3116.5147\text{ s}}
	\end{equation}
	
	In condensed matter configurations, substrate density engineering can utilize this suppressed noise floor to isolate processing states from ambient thermal noise at room temperature ($293.15\text{ K}$). Under edge-reflected packet filter configurations, localized stochastic pressure gradients vanish cleanly ($\max c_s^2 = 9.2738 \times 10^{-10}$), minimizing environmental phonon collisions and preserving a local coherence fidelity plateau of approximately $78.97\%$.
	
	\begin{thebibliography}{99}
		
		\bibitem{longo2026}
		P. A. Longo, 
		\textit{A Six-Dimensional Torsional Action Principle and its Explicit Resolution of the Cosmological Hubble Tension}, 
		Zenodo Preprint, \url{https://doi.org} (2026). 
		
		\bibitem{Khoury2004}
		J. Khoury and A. Weltman,
		\textit{Chameleon fields: Awaits gravity mediated by an effective scalar field},
		Phys. Rev. D \textbf{69}, 044026 (2004).
		
		\bibitem{Magueijo2003}
		J. Magueijo,
		\textit{New varying speed of light theories},
		Rept. Prog. Phys. \textbf{66}, 2025 (2003).
		
		\bibitem{Riess2022}
		A. G. Riess et al.,
		\textit{A Comprehensive Measurement of the Local Value of the Hubble Constant with $1\sigma$ Accuracy from the Hubble Space Telescope and the SH0ES Team},
		Astrophys. J. Lett. \textbf{934}, L7 (2022).
		
		\bibitem{Planck2020}
		Planck Collaboration,
		\textit{Planck 2018 results. VI. Cosmological parameters},
		Astron. Astrophys. \textbf{641}, A6 (2020).
		
	\end{thebibliography}
\end{document}
